% 由 画表图_ai.jl 自动生成，请勿手改；改数据请改 表图数据.xlsx 后重跑
\begin{table}[htbp]
    \centering
    \caption{四组网格配置、四档网格密度下 $w$、$\varphi$ 与 $Q$ 的相对 $L_2$ 误差（$\log_{10}$）}
    \label{tab:conv4}
    \small
\renewcommand{\arraystretch}{1.25}
\setlength{\tabcolsep}{6pt}
\resizebox{\textwidth}{!}{%
\begin{tabular}{@{}lcccccccccccc@{}}
\toprule[1.0pt]
 & \multicolumn{3}{c}{tri3 均匀网格} & \multicolumn{3}{c}{tri3 W 非均匀网格} & \multicolumn{3}{c}{quad4 均匀网格} & \multicolumn{3}{c}{quad4 W 非均匀网格} \\
\cmidrule(lr){2-4}\cmidrule(lr){5-7}\cmidrule(lr){8-10}\cmidrule(lr){11-13}
$n_q$ & $w$ & $\varphi$ & $Q$ & $w$ & $\varphi$ & $Q$ & $w$ & $\varphi$ & $Q$ & $w$ & $\varphi$ & $Q$ \\
\midrule[0.5pt]
81 & $-1.182529$ & $-1.023714$ & $-0.898027$ & $-1.106342$ & $-0.971451$ & $-0.337787$ & $-1.223624$ & $-1.110703$ & $-0.837836$ & $-1.082464$ & $-1.069857$ & $-0.562672$ \\
289 & $-1.782627$ & $-1.577969$ & $-1.543729$ & $-1.656352$ & $-1.492693$ & $-0.310635$ & $-1.935099$ & $-1.752441$ & $-1.460734$ & $-1.712445$ & $-1.491926$ & $-0.167035$ \\
1089 & $-2.398376$ & $-2.154196$ & $-2.116972$ & $-2.020997$ & $-1.866298$ & $0.003165$ & $-2.494315$ & $-2.332677$ & $-2.013936$ & $-2.167660$ & $-1.925892$ & $-0.045204$ \\
4225 & $-3.007518$ & $-2.740824$ & $-2.554353$ & $-2.216214$ & $-2.139113$ & $0.254666$ & $-3.038233$ & $-2.908195$ & $-2.323857$ & $-2.137562$ & $-2.044939$ & $0.265515$ \\
\bottomrule[1.0pt]
\end{tabular}}%

\par\vspace{3pt}
{\footnotesize\raggedright
$w$ 为无网格 RK 场；$\varphi$ 与 $q$ 为有限元场（同网格、同单元）\\
$n_q$ 为 $q$ 与 $\varphi$ 的节点数；表中数值为 $\log_{10}$ 相对 $L_2$ 误差\\
非均匀网格组（W 网格）的支撑尺寸按单元面积逐节点确定
\par}

\end{table}
